\documentclass{article}
\usepackage{amsmath, amssymb, amsthm}
\usepackage{hyperref}
\usepackage{geometry}
\geometry{margin=1in}

\title{Erd\H{o}s Problem \#1107}
\date{Last updated: 2025-11-17}
\author{Source: erdosproblems.com}

\begin{document}
\maketitle

\noindent\textbf{Status:} OPEN \quad \textbf{No prize}

\noindent\textbf{Tags:} number theory, powerful

\medskip
\noindent\textbf{Problem Statement:}

Let $r\geq 2$. A number $n$ is $r$-powerful if for every prime $p$ which divides $n$ we have $p^r\mid n$. Is every large integer the sum of at most $r+1$ many $r$-powerful numbers?

\bigskip
\noindent\textbf{Remarks:}

Given in the 1986 Oberwolfach problem book as a problem of Erd\H{o}s and Ivi\'{c}.

This is true when $r=2$, as proved by Heath-Brown \cite{He88} (see [941]). 

See [940] for the problem of which integers are the sum of at most $r$ many $r$-powerful numbers.

\bigskip
\noindent\textbf{References:}

[He88] Heath-Brown, D. R., Ternary quadratic forms and sums of three square-full numbers. (1988), 137--163.

\bigskip
\noindent\small{Source: \url{https://www.erdosproblems.com/1107}}
\end{document}
